\vspace*{0.5cm}
\noindent\begin{tabularx}{\textwidth}{@{}>{\bfseries}l@{}
@{\hspace{1cm}}l@{} @{\hspace{0.5cm}}l@{}}
\textthai{หัวข้อวิทยานิพนธ์}		  & \multicolumn{2}{@{}X@{}}{\begin{tabular}[t]{@{}X@{}}\textthai{ชื่อวิทยานิพนธ์ภาษาไทยของคุณบรรทัดที่ 1} \\ \textthai{บรรทัดที่ 2 ถ้าชื่อยาว}\end{tabular}}\vspace{11pt}\\
\textthai{ผู้เขียน}				& \multicolumn{2}{@{}l@{}}{\printNamePrefixInThai\printAuthorInThai}\vspace{11pt}\\
\textthai{ปริญญา}				 & \multicolumn{2}{@{}l@{}}{\printDegreeInThai(\printProgramInThai)}\vspace{11pt}\\
\textthai{คณะกรรมการที่ปรึกษา}		& \printAdvisorInThai			& \textthai{อาจารย์ที่ปรึกษาหลัก}\\
						& \printCoadvisorAInThai	& \textthai{อาจารย์ที่ปรึกษาร่วม}\\
\end{tabularx}


\begin{ThaiAbstract}

% Write your Thai abstract here.
% Wrap all Thai text in \textthai{...}
%
% โครงสร้างที่แนะนำ:
% 1. ความเป็นมาและแรงจูงใจ (2-3 ประโยค)
% 2. วัตถุประสงค์การวิจัย (1-2 ประโยค)
% 3. ระเบียบวิธีวิจัย (3-5 ประโยค)
% 4. ผลการวิจัยที่สำคัญ (3-5 ประโยค)
% 5. บทสรุป (1-2 ประโยค)

\textthai{บทคัดย่อภาษาไทยของคุณอยู่ที่นี่}

\end{ThaiAbstract}
